\documentclass{article}

\title{中文排版：从断行到成书}
\subtitle{soma-typeset 首发示例}
\author{zhanghuanhuan303}
\date{2026-11-20}

\begin{document}

\section{为什么是全局优化}

排版引擎要做的第一件事不是画字，而是决定每一行在哪里结束。这个决定一旦做错，后面所有的字号、行距和页边距都只是在修饰一个错误。中文排版尤其如此：汉字等宽，标点占满一格，一旦允许标点出现在行首，读者的视线就会在换行处被切断。因此，禁则、挤压和悬挂不是装饰性的细节，而是断行模型的一部分。

\begin{figure}
\centering
\includegraphics[width=0.7\linewidth]{fig/pipeline.pdf}
\caption{渲染管线：语法层、统一 IR 与排版引擎}
\end{figure}

\subsection{两条设计原则}

我们在设计代价模型时坚持两条原则。第一，每一类惩罚都必须有独立的权重，能够被单独观察和标定；如果两项代价在公式里相乘，调参就会退化成黑箱。第二，所有几何量都必须是整数，跨平台的结果必须逐字节一致，否则一次浏览器升级就可能让同一份文档换一个断点。这两条原则看起来保守，却让后续每一次优化都有可回退的基线\footnote{整数字面量以 1/4096 em 为单位，除法一律向零截断。}。

\subsection{断行不是对齐}

一个常见的误解是：只要把行末对齐，版面就整齐了。事实并非如此。行与行之间的松紧差异过大时，页面会出现明显的明暗条纹；段落与段落之间的空隙如果全靠固定间距，页面底部就会留下无法解释的空白。可拉伸的垂直胶水正是为此存在：它让页面在多一行与少一行之间有一个连续的中间状态，而不是非此即彼。

\section{度量从哪里来}

我们把字体度量离线提取成一张整数表，运行时只查表，不解析字体。表里记录每个字形的宽度、字面的上伸下伸，以及空格的宽度。这样做有三个直接后果：构建不依赖字体文件；native 与 WebAssembly 使用同一张表；换字体就是换一张表，哈希随之改变，结果仍然可复核。

\subsection{一张表能做与不能做的事}

需要说明的是，这张表只包含数值，不包含字形轮廓。它不是字体的替代品，而是排版决策所需的度量摘要。真正的字形渲染仍然交给浏览器或打印后端，引擎只负责告诉它们“这个字放在这里，占这么宽”。

数学符号是度量表里的一个特例。求和号、积分号、根号的宽度与正文完全不同，如果借用正文字体，根号会与它的横线分离，求和号的上下限也会挤在一起。因此我们为公式单独准备了一面字体度量，第二遍排版时只换这一面，行内公式的宽度与行盒高度就都正确了。

表格是另一个特例。单元格的宽度由列数决定，而列宽又依赖内容的自然宽度；当单元格里的文字要换行时，列宽与行数是同时确定的。我们选择先量化、再排版：先按列数平分版心，再做行内断行，最后把单元格的高度取行数的最大值。这样每一行都能整除，续页的表头也能精确重现。

长文档里，这些数字会累积。一张表跨三页，表头要在每一页重画，行高要保持一致，跨行的单元不能被切开；任何一处偏差都会在第十页变成明显的错位。所以我们宁可让表格的模型朴素一些，也不允许它在分页时失去可解释性。

\section{垂直方向与脚注}

垂直方向的度量比水平方向简单，因为行高是固定的、可加的，页面一旦排满就可以停下来。这个性质还带来一个额外的好处：分页的动态规划可以精确剪枝，长文档不会因此变慢。真正困难的是脚注：脚注的高度会反过来影响它所在的页面，必须和正文一起求解，而不能先排正文再补脚注。

我们把脚注当作所在页面的固定开销：先估计脚注区的高度，再在剩下的空间里排正文，排不下就把最后一行连同它的脚注一起推到下一页。迭代最多八轮，超过轮数就保留最后一个稳定结果，并在页面上标记出来。这样做不优雅，但每一步都可以复算。

行高与基线网格之间的关系也值得说明。网格是对齐的工具，不是必须；当模板指定了网格时，每一行的行高会被量化到网格的整数倍，行与行之间的基线就会严格等距。代价是页面底部的余量可能多出半行，这部分余量无法再用胶水吸收，我们会把它记在页面的空白里。

度量表的哈希会进入页面结果。这意味着换字体、换字号表、甚至重新生成表，都会改变结果哈希；反过来，只要哈希一致，两次排版的每一行坐标就必须完全一致。这条约束是我们所有跨端承诺的基础，也是回归测试里最便宜的一道门。

有人会问：既然 TeX 已经把断行做了四十年，为什么不直接用它？原因是我们要的是一个可以被嵌入、被测试、被逐字节复核的核心。TeX 的断行与它的宏语言、字体系统、页面模型绑在一起，抽出来的成本远高于重写；而重写之后，每一条代价、每一个阈值都能写进参数表，出了问题可以回退到某一个版本，而不是回退到某一次运气。

整数与浮点是另一个容易低估的差异。浮点运算在不同平台、不同编译器优化下可能给出不同的最后一位，而断行是离散决策：比例差一点点，断点就可能换一个位置。我们把所有几何量都换成整数，除法一律向零截断，代价是精度只有 1/4096 em，好处是同一份文档在任何目标上都会得到同一份结果。

可复核还包括“同样的输入给出同样的哈希”。文档哈希、模板哈希、度量哈希、参数哈希、引擎版本，五个值一起决定页面结果。任何一个变了，哈希就变；哈希不变，两次排版的每一行坐标就必须一致。这条约束让回归测试变得便宜：比较一个十六进制数就够了。

\section{标点、边界与行边裁剪}

中文标点的处理可以拆成三件事：哪些位置不能断，哪些位置可以压，哪些标点可以伸出行框。

\begin{itemize}
\item 禁则：行首不放收尾标点，行尾不留起始标点。这是硬约束，宁可让一行略紧，也不让标点落在不该落的地方。
\item 挤压：行尾的句号、逗号、右引号可以压掉一部分字面空白；行首的左引号、左括号同理。
\item 悬挂：允许收尾标点越过行尾一小段，视觉上反而更齐。
\end{itemize}

这三件事都不应该改动断行算法本身。禁则是断点属性，挤压是行边裁剪，悬挂是排版完成后的坐标修正。把它们混进代价函数，权重就再也解释不清。

中英混排还有一个额外问题：汉字与拉丁字母之间的间隙。理想情况下，这个间隙随行内调整按比例伸缩；实际实现里，它只能是一个近似值。我们用行边偏差作为观测量，把近似值与真实值都记录下来，偏差超过阈值就报警，而不是假装它不存在。

\section{长文、目录与章节页}

短文档靠肉眼，长文档靠机制。三十页之后，孤行、寡行、页面底部的空白、连续两页的松紧差异都会累积成明显的节奏问题。我们要求每一页的填充都可以解释：要么正文排满，要么脚注占位，要么是章节起始的刻意留白。

目录是长文的第一个考验。目录本身的页数会改变正文页码，而页码又是目录的内容，这是一个不动点问题。我们的做法是最多迭代四遍：用上一遍的页码排版，从分页结果反解标题所在页，收敛就停；不收敛时在结果里标记为不稳定，而不是假装已经收敛。

章节页则是版面的呼吸。章首不排页眉，让标题独占版心上方；正文从下一页开始，页眉恢复为当前章名。这个规则看起来简单，却要求分页器知道哪一页是章节的起点，而不是让作者手写分页符。

\section{附录：公式与表格}

行内公式参与断行，行间公式独占一块垂直空间：

\[ \left( \begin{aligned} \text{行宽} &= W - 2m \\ \text{压缩量} &= \frac{s}{S} \cdot D \end{aligned} \right) \]

当行内调整不足以吸收差额时，引擎会在应急阶段放宽容忍度，并在结果中逐行标记，不让一次降级悄悄发生。

\begin{table}
\caption{三个阶段与可见后果}
\begin{tabular}{lll}
\textbf{阶段} & \textbf{触发条件} & \textbf{可见后果} \\ \midrule
正常 & 比例在容忍度内 & 无 \\
应急 & 比例超出容忍度 & 行标记为 emergency \\
救援 & 硬禁则无解 & 允许跨越禁则并告警 \\
\end{tabular}
\end{table}

\begin{itemize}
\item 正常阶段：断点由动态规划给出，代价可解释。
\item 应急阶段：放宽比例上限，行仍保持整数坐标。
\item 救援阶段：只在无解时触发，结果必须能被审计。
\end{itemize}

\end{document}
